% GENERATED FILE. Edit canonical lessons, then run npm run generate:books.
% canonical-source-hash: fnv1a64:77a80a0cf6093bfd
% canonical-lessons: TA-C508-utiyam, TA-C508-vituppu, TA-C508-kututal-neram, TA-C508-kanakkalar, TA-C508-eluttar

\chapter{Pay and Leave}
\label{ch:ta-pay-and-leave}
\hlchaptermodality{508}

\begin{chapteropening}
\textbf{By the end of this chapter:} \emph{I can talk about wages, leave and overtime.}

Complete the last lesson of chapter 508: I can talk about wages, leave and overtime.
\end{chapteropening}

\section[\texorpdfstring{\ta{ஊதியம்} (ūtiyam)}{ஊதியம் (ūtiyam)}]{\ta{ஊதியம்} (ūtiyam) — wages; pay}

\label{lesson:TA-C508-utiyam}



\begin{quote}
Before the new one: say the Tamil for a file, then the Tamil for to supervise.
\end{quote}

\subsection*{You'll want to know: \ta{ஊதியம்}}
\textbf{\ta{ஊதியம்}} — \emph{ūtiyam} — \textquotedblleft{}wages; pay\textquotedblright{}.

\begin{quote}
\ta{இந்த} \ta{மாதம்} \ta{ஊதியம்} \ta{தாமதமாக} \ta{வந்தது}.
\end{quote}

\begin{quote}
\emph{inta mātam ūtiyam tāmatamāka vantatu.}
\end{quote}

\begin{quote}
This month the pay came late.
\end{quote}

A formal word beside \ta{சம்பளம்} (\emph{campaḷam}); it is the second half of \ta{ஓய்வூதியம்} (\emph{ōyvūtiyam}), a pension, the pay of rest.

\begin{cousinweb}
The origin of \textbf{\ta{ஊதியம்}} (\emph{ūtiyam}) is not settled: DEDR 609, under \ta{உதவு} (\emph{utavu}), \textquotedblleft{}to help, to be of use\textquotedblright{}, lists \ta{ஊதியம்}, \textquotedblleft{}profit, gain\textquotedblright{}, but marks the link with a question mark.
\end{cousinweb}

\begin{grammarlens}[title={What you've built}]
One more word for describing what you lived and felt.
\end{grammarlens}

\subsection*{Your turn}
\begin{itemize}
\raggedright
  \item \emph{Say it:} \emph{ūtiyam}
  \item \emph{Say it:} \emph{ūtiyam}, once more
  \item \emph{Say it:} \emph{inta mātam ūtiyam tāmatamāka vantatu.} — This month the pay came late.
  \item \emph{Recall:} say the Tamil for a file, then the Tamil for to supervise, then say \emph{ūtiyam} again
\end{itemize}

\begin{tcolorbox}[breakable,colback=teal!4,colframe=teal!35!black,title={Before you move on}]
What does \ta{ஊதியம்} mean? (\textquotedblleft{}Wages; pay\textquotedblright{}.) Say it once more.
\end{tcolorbox}
\section[\texorpdfstring{\ta{விடுப்பு} (viṭuppu)}{விடுப்பு (viṭuppu)}]{\ta{விடுப்பு} (viṭuppu) — leave from work}

\label{lesson:TA-C508-vituppu}



\begin{quote}
Before the new one: say the Tamil for to supervise, then the Tamil for wages.
\end{quote}

\subsection*{You'll want to know: \ta{விடுப்பு}}
\textbf{\ta{விடுப்பு}} — \emph{viṭuppu} — \textquotedblleft{}leave from work\textquotedblright{}.

\begin{quote}
\ta{நான்} \ta{நாளை} \ta{ஒரு} \ta{நாள்} \ta{விடுப்பு} \ta{எடுக்கிறேன்}.
\end{quote}

\begin{quote}
\emph{nāṉ nāḷai oru nāḷ viṭuppu eṭukkiṟēṉ.}
\end{quote}

\begin{quote}
I am taking a day of leave tomorrow.
\end{quote}

One takes leave with \ta{எடு} (\emph{eṭu}): \ta{விடுப்பு} \ta{எடு}. A holiday for everyone is \ta{விடுமுறை} (\emph{viṭumuṟai}).

\begin{cousinweb}
\textbf{\ta{விடுப்பு}} (\emph{viṭuppu}) is native Dravidian, a noun made from \ta{விடு} (\emph{viṭu}), \textquotedblleft{}to leave, to let go, to release\textquotedblright{} (DEDR 5393), with the ending -\ta{ப்பு}.
\end{cousinweb}

\begin{grammarlens}[title={What you've built}]
One more word for describing what you lived and felt.
\end{grammarlens}

\subsection*{Your turn}
\begin{itemize}
\raggedright
  \item \emph{Say it:} \emph{viṭuppu}
  \item \emph{Say it:} \emph{viṭuppu}, once more
  \item \emph{Say it:} \emph{nāṉ nāḷai oru nāḷ viṭuppu eṭukkiṟēṉ.} — I am taking a day of leave tomorrow.
  \item \emph{Recall:} say the Tamil for to supervise, then the Tamil for wages, then say \emph{viṭuppu} again
\end{itemize}

\begin{tcolorbox}[breakable,colback=teal!4,colframe=teal!35!black,title={Before you move on}]
What does \ta{விடுப்பு} mean? (\textquotedblleft{}Leave from work\textquotedblright{}.) Say it once more.
\end{tcolorbox}
\section[\texorpdfstring{\ta{கூடுதல்} \ta{நேரம்} (kūṭutal nēram)}{கூடுதல் நேரம் (kūṭutal nēram)}]{\ta{கூடுதல்} \ta{நேரம்} (kūṭutal nēram) — overtime}

\label{lesson:TA-C508-kututal-neram}



\begin{quote}
Before the new one: say the Tamil for wages, then the Tamil for leave from work.
\end{quote}

\subsection*{You'll want to know: \ta{கூடுதல்} \ta{நேரம்}}
\textbf{\ta{கூடுதல்} \ta{நேரம்}} — \emph{kūṭutal nēram} — \textquotedblleft{}overtime\textquotedblright{}.

\begin{quote}
\ta{இந்த} \ta{வாரம்} \ta{நான்} \ta{கூடுதல்} \ta{நேரம்} \ta{வேலை} \ta{செய்தேன்}.
\end{quote}

\begin{quote}
\emph{inta vāram nāṉ kūṭutal nēram vēlai ceytēṉ.}
\end{quote}

\begin{quote}
This week I worked overtime.
\end{quote}

Word for word \textquotedblleft{}extra time\textquotedblright{}: \ta{கூடுதல்} \textquotedblleft{}extra, more\textquotedblright{} + \ta{நேரம்} \textquotedblleft{}time\textquotedblright{}.

\begin{cousinweb}
\textbf{\ta{கூடுதல்} \ta{நேரம்}} (\emph{kūṭutal nēram}) is a native phrase: \ta{கூடுதல்} (\emph{kūṭutal}), \textquotedblleft{}extra\textquotedblright{}, from \ta{கூடு} (\emph{kūṭu}), \textquotedblleft{}to come together, to increase\textquotedblright{} (DEDR 1882), and \ta{நேரம்} (\emph{nēram}), \textquotedblleft{}time\textquotedblright{} (DEDR 3774).
\end{cousinweb}

\begin{grammarlens}[title={What you've built}]
One more word for describing what you lived and felt.
\end{grammarlens}

\subsection*{Your turn}
\begin{itemize}
\raggedright
  \item \emph{Say it:} \emph{kūṭutal nēram}
  \item \emph{Say it:} \emph{kūṭutal nēram}, once more
  \item \emph{Say it:} \emph{inta vāram nāṉ kūṭutal nēram vēlai ceytēṉ.} — This week I worked overtime.
  \item \emph{Recall:} say the Tamil for wages, then the Tamil for leave from work, then say \emph{kūṭutal nēram} again
\end{itemize}

\begin{tcolorbox}[breakable,colback=teal!4,colframe=teal!35!black,title={Before you move on}]
What does \ta{கூடுதல்} \ta{நேரம்} mean? (\textquotedblleft{}Overtime\textquotedblright{}.) Say it once more.
\end{tcolorbox}
\section[\texorpdfstring{\ta{கணக்காளர்} (kaṇakkāḷar)}{கணக்காளர் (kaṇakkāḷar)}]{\ta{கணக்காளர்} (kaṇakkāḷar) — an accountant}

\label{lesson:TA-C508-kanakkalar}



\begin{quote}
Before the new one: say the Tamil for leave from work, then the Tamil for overtime.
\end{quote}

\subsection*{You'll want to know: \ta{கணக்காளர்}}
\textbf{\ta{கணக்காளர்}} — \emph{kaṇakkāḷar} — \textquotedblleft{}an accountant\textquotedblright{}.

\begin{quote}
\ta{கணக்காளர்} \ta{எல்லா} \ta{ரசீதுகளையும்} \ta{சரிபார்த்தார்}.
\end{quote}

\begin{quote}
\emph{kaṇakkāḷar ellā racītukaḷaiyum caripārttār.}
\end{quote}

\begin{quote}
The accountant checked all the receipts.
\end{quote}

\ta{கணக்கு} \textquotedblleft{}an account, a sum\textquotedblright{} + -\ta{ஆளர்} \textquotedblleft{}one who …\textquotedblright{}. The plural is \ta{கணக்காளர்கள்}.

\begin{cousinweb}
\textbf{\ta{கணக்காளர்}} (\emph{kaṇakkāḷar}) is \ta{கணக்கு} (\emph{kaṇakku}), \textquotedblleft{}an account, a sum\textquotedblright{}, with -\ta{ஆளர்}, \textquotedblleft{}one who does it\textquotedblright{}. The origin of \ta{கணக்கு} is not settled.
\end{cousinweb}

\begin{grammarlens}[title={What you've built}]
One more word for describing what you lived and felt.
\end{grammarlens}

\subsection*{Your turn}
\begin{itemize}
\raggedright
  \item \emph{Say it:} \emph{kaṇakkāḷar}
  \item \emph{Say it:} \emph{kaṇakkāḷar}, once more
  \item \emph{Say it:} \emph{kaṇakkāḷar ellā racītukaḷaiyum caripārttār.} — The accountant checked all the receipts.
  \item \emph{Recall:} say the Tamil for leave from work, then the Tamil for overtime, then say \emph{kaṇakkāḷar} again
\end{itemize}

\begin{tcolorbox}[breakable,colback=teal!4,colframe=teal!35!black,title={Before you move on}]
What does \ta{கணக்காளர்} mean? (\textquotedblleft{}An accountant\textquotedblright{}.) Say it once more.
\end{tcolorbox}
\section[\texorpdfstring{\ta{எழுத்தர்} (eḻuttar)}{எழுத்தர் (eḻuttar)}]{\ta{எழுத்தர்} (eḻuttar) — a clerk}

\label{lesson:TA-C508-eluttar}



\begin{quote}
Before the new one: say the Tamil for overtime, then the Tamil for an accountant.
\end{quote}

\subsection*{You'll want to know: \ta{எழுத்தர்}}
\textbf{\ta{எழுத்தர்}} — \emph{eḻuttar} — \textquotedblleft{}a clerk\textquotedblright{}.

\begin{quote}
\ta{அலுவலகத்தில்} \ta{ஒரு} \ta{எழுத்தர்} \ta{என்} \ta{படிவத்தை} \ta{வாங்கினார்}.
\end{quote}

\begin{quote}
\emph{aluvalakattil oru eḻuttar eṉ paṭivattai vāṅkiṉār.}
\end{quote}

\begin{quote}
At the office a clerk took my form.
\end{quote}

\ta{எழுத்து} \textquotedblleft{}writing\textquotedblright{} + -\ta{அர்} \textquotedblleft{}person\textquotedblright{}: the one who writes the papers.

\begin{cousinweb}
\textbf{\ta{எழுத்தர்}} (\emph{eḻuttar}) is native Dravidian: \ta{எழுத்து} (\emph{eḻuttu}), \textquotedblleft{}writing\textquotedblright{}, from \ta{எழுது} (\emph{eḻutu}), \textquotedblleft{}to write\textquotedblright{} (DEDR 853), with the personal ending -\ta{அர்}.
\end{cousinweb}

\begin{grammarlens}[title={What you've built}]
One more word for describing what you lived and felt.
\end{grammarlens}

\subsection*{Your turn}
\begin{itemize}
\raggedright
  \item \emph{Say it:} \emph{eḻuttar}
  \item \emph{Say it:} \emph{eḻuttar}, once more
  \item \emph{Say it:} \emph{aluvalakattil oru eḻuttar eṉ paṭivattai vāṅkiṉār.} — At the office a clerk took my form.
  \item \emph{Recall:} say the Tamil for overtime, then the Tamil for an accountant, then say \emph{eḻuttar} again
\end{itemize}

\begin{tcolorbox}[breakable,colback=teal!4,colframe=teal!35!black,title={Before you move on}]
What does \ta{எழுத்தர்} mean? (\textquotedblleft{}A clerk\textquotedblright{}.) Say it once more.
\end{tcolorbox}
